% sections/intro.tex -- every number here is a macro from tables/numbers.tex.
\section{Introduction}
\label{sec:intro}

A learned surrogate for a family of PDE solves is usually a single point in the
accuracy--cost plane. A neural operator \citep{Li2020FNO,Lu2021DeepONet,
Kovachki2023NeuralOperator} evaluates in one forward pass at one accuracy;
getting a more accurate or a cheaper answer means training another model. A
projection-based reduced-order model (ROM) has the opposite character: because
it solves the discretised equations in a small number of coordinates, it
inherits the solver's controls, and a classical linear ROM exposes an obvious
family through its basis rank. A \emph{nonlinear-manifold} ROM
\citep{LeeCarlberg2020,KimChoi2022NMROM} trades that rank for a learned
decoder whose few coordinates reach further than the same number of linear
modes, and in doing so appears to give the family up: the decoder is fixed once
trained, and the solver's remaining knobs move cost far more than they move
accuracy.

This paper is about recovering the family without retraining, and about being
precise on what it is worth. We work with a decoder that writes the solution as
a fixed spatial \emph{bank} $\bank$ of $R$ functions times a small neural
\emph{head} $\head$ of $k$ latent coordinates, and at inference we let the
solver add $q$ fixed linear directions on top of the head's output,
$u=\bank(\head(\latent)+\corr y)$. The rank $q$ dials the reachable set from
the head's image ($q=0$) to the bank's whole span ($q=R$), where the model is a
linear ROM and needs no nonlinear iteration at all. The paper has the same
three-part skeleton as our earlier attempt at this claim, re-founded on this
decoder and on evidence we now trust, and it reports the losses as losses.

\paragraph{Contributions.}
\begin{enumerate}
\item \textbf{A tunable NM-ROM from one trained decoder.} With the test count
  $M$ held fixed so that $q$ is the only control, the Burgers $256^2$ ladder is
  monotone, every rung converges, and its four rungs span
  $\nQxmErrSpan\times$ in error and $\nQxmCostSpan\times$ in cost inside one
  allocation (\S\ref{sec:exp:qxm}); this meets a bar fixed before any run. The
  cost side has two further levers, offline-certified quadrature rules
  ($\nEqtopCostRatioMin$--$\nEqtopCostRatioMax\times$ the dense cost at
  matched error, provisional) and the stopping tolerance
  (\S\ref{sec:exp:panel}). At matched latent dimension the family beats the
  classical linear ROM (\S\ref{sec:exp:head}); it does not beat a tuned
  full-order solver or a well-trained neural operator at $256^2$, and we say
  so first (\S\ref{sec:exp:panel}, \S\ref{sec:exp:operators}).
\item \textbf{The architecture that makes it possible.} A bank that evaluates
  per node, so the cached reduced cost is flat while the mesh grows
  $\nMeshBurgersUnknownGrowth\times$ (\S\ref{sec:exp:mesh}); a linear skip in
  the head, so the cold solve starts reliably; and the bank/head/correction
  split, which is what lets a run-time rank interpolate between the head's
  compression and the bank's floor (\S\ref{sec:method}).
\item \textbf{The structural condition.} Where the weak residual is nonlinear in
  the coefficients (Burgers), the corrections cost a growing nonlinear solve and
  the ladder is a genuine trade. Where it is linear (Poisson, heat, waves) the
  corrections are eliminated exactly, the top rung is the linear model, and it
  is also the cheapest point: the family collapses, and on waves the learned
  bank is itself worse than a POD basis of the same rank
  (\S\ref{sec:exp:linear}). This says when the nonlinear manifold is worth
  having.
\item \textbf{Two methodological findings.} Every result is reported in a
  three-layer decomposition---bank floor, best-found, solved---so that each
  lever's effect can be assigned to the layer it acts on
  (\S\ref{sec:method:layers}); and quadrature rules must be certified on states
  the solver actually reaches, never by their fitting residual
  (\S\ref{sec:exp:eq}).
\end{enumerate}

\paragraph{What we do not claim.} No speedup over an efficient full-order
solver in 2D: at $256^2$ none of the $\nPanelReducedCount$ reduced-order
subjects in the same-allocation panel is non-dominated once the full-order
controls are included. No accuracy superiority over neural operators: a U-Net
and a Transolver trained on the same data beat the ROM on the matched cohort,
and only the FNO does not. No frontier over classical POD at every rank:
POD-512 beats the best dense rung on the evolved metric. No cost ratio across
jobs, ever. No convergence theory. And every number is single-seed,
development-cohort evidence until the seed and sealed-cohort runs land
(\S\ref{sec:limitations}).

\paragraph{The reviewers' asks.} This is a resubmission, and the three reviews
of the earlier version asked for specific things: one checkpoint per curve and
a stated protocol (fxe8); competitive full-order solvers, precisely specified
PDEs and harder domains (GwrW); tuned baselines on shared data, seeds, and a
case where the nonlinear manifold is the enabling ingredient (5mgh). Each ask is
mapped to the section that answers it, or to the limitation that concedes it,
in Appendix~\ref{app:reviewers}.
